% Shared preamble of the technical report, French (rapport.tex) and English (report-en.tex).
% The main file sets \reportlang to fr or en before % Shared preamble of the technical report, French (rapport.tex) and English (report-en.tex).
% The main file sets \reportlang to fr or en before % Shared preamble of the technical report, French (rapport.tex) and English (report-en.tex).
% The main file sets \reportlang to fr or en before % Shared preamble of the technical report, French (rapport.tex) and English (report-en.tex).
% The main file sets \reportlang to fr or en before \input{preamble}.
% LuaLaTeX + biber. Fonts: Latin Modern for text, headings and code, Latin Modern Math for formulas;
% the figures are drawn in the same Latin Modern.
\documentclass[11pt,a4paper,oneside]{report}
\usepackage[a4paper,top=2.6cm,bottom=2.8cm,inner=2.7cm,outer=2.7cm,headheight=14pt]{geometry}

% The LaTeX fonts throughout: Latin Modern for text, headings and code, Latin Modern Math for formulas
\usepackage{fontspec}
% The ffi and ffl ligatures are set as ff + i and ff + l: Ghostscript (build.sh) loses their ToUnicode entry,
% and pdftotext, search and copy-paste then read « ofÏciel ». The ff, fi and fl ligatures survive it.
\directlua{fonts.handlers.otf.addfeature{name="noffi",type="multiple",data={["f_f_i"]={"f_f","i"},["f_f_l"]={"f_f","l"}}}}
\defaultfontfeatures{RawFeature=+noffi}
\setmainfont{Latin Modern Roman}
\setsansfont{Latin Modern Sans}
\setmonofont{Latin Modern Mono}
\newcommand{\headfont}{\rmfamily}

\def\langfr{fr}
\newcommand{\iffr}[2]{\ifx\reportlang\langfr #1\else #2\fi}
\ifx\reportlang\langfr \usepackage[french]{babel}\frenchsetup{StandardItemLabels=true,ItemLabels=\textbullet}\else \usepackage[british]{babel}\fi
\usepackage{csquotes}
\usepackage{microtype}
\usepackage{amsmath,amssymb}
\usepackage{unicode-math}
\setmathfont{Latin Modern Math}
\usepackage{siunitx}
\ifx\reportlang\langfr \sisetup{locale=FR,group-minimum-digits=4,mode=text}\else \sisetup{locale=UK,group-separator={,},group-minimum-digits=4,mode=text}\fi

% Colours of the figures (article/figures/draw.py)
\usepackage[table]{xcolor}
\definecolor{ink}{HTML}{15171C}
\definecolor{ink2}{HTML}{4A4F58}
\definecolor{muted}{HTML}{868C96}
\definecolor{rulec}{HTML}{C9CDD3}
\definecolor{panel}{HTML}{F5F6F8}
\definecolor{accent}{HTML}{D6453D}
\definecolor{okgreen}{HTML}{2F9E5B}
\definecolor{sacblue}{HTML}{2F6FD6}
\definecolor{crossqorange}{HTML}{EE8320}
\definecolor{ppotorch}{HTML}{C23B75}
\definecolor{ppojax}{HTML}{109C86}
\definecolor{link}{HTML}{2F6FD6}

\usepackage{graphicx}
\graphicspath{{figures/}}
\usepackage{booktabs,longtable,array,tabularx,multirow,threeparttable}
\usepackage{makecell}
\renewcommand{\arraystretch}{1.15}
\usepackage[font={small},labelfont={bf},labelsep=period,justification=raggedright,singlelinecheck=false]{caption}
% Captions above tables (the generated tables put them first), below figures
\captionsetup[table]{position=top,skip=6pt}
\captionsetup[figure]{position=bottom,skip=8pt}
% Floats: fill text pages before making a float page, never a float page with one small figure and a blank half
\usepackage{placeins}
\usepackage{flafter}% a float never comes before the text that calls it
% A long table never starts with its caption and header alone at the foot of a page
\usepackage{needspace}
\BeforeBeginEnvironment{longtable}{\par\Needspace*{14\baselineskip}}
\renewcommand{\topfraction}{0.88}
\renewcommand{\bottomfraction}{0.6}
\renewcommand{\textfraction}{0.1}
\renewcommand{\floatpagefraction}{0.72}
\setcounter{topnumber}{3}
\setcounter{bottomnumber}{2}
\setcounter{totalnumber}{4}
\setlength{\floatsep}{14pt plus 4pt minus 2pt}
\setlength{\textfloatsep}{18pt plus 4pt minus 4pt}
\setlength{\intextsep}{14pt plus 4pt minus 2pt}
\makeatletter
\setlength{\@fptop}{0pt}% float pages fill from the top
\setlength{\@fpsep}{14pt plus 1fil}
\setlength{\@fpbot}{0pt plus 1fil}
\makeatother
\usepackage{subcaption}
\usepackage{enumitem}
\setlist{itemsep=2pt,topsep=4pt}
\usepackage{listings}
\lstset{basicstyle=\ttfamily\small,columns=fullflexible,keepspaces=true,breaklines=true,
  frame=l,framerule=1.2pt,rulecolor=\color{rulec},xleftmargin=1.2em,backgroundcolor=\color{panel},
  showstringspaces=false,commentstyle=\color{muted},keywordstyle=\color{ink}\bfseries}
\usepackage[most]{tcolorbox}
\newtcolorbox{keybox}[1][]{enhanced,colback=panel,colframe=accent,boxrule=0pt,leftrule=3pt,arc=0pt,
  left=10pt,right=10pt,top=6pt,bottom=6pt,fontupper=\small,#1}
\newtcolorbox{notebox}[1][]{enhanced,colback=white,colframe=rulec,boxrule=0.6pt,arc=2pt,
  left=8pt,right=8pt,top=5pt,bottom=5pt,fontupper=\small,#1}

% Headings
\usepackage{titlesec}
\titleformat{\chapter}[display]{\headfont\color{ink}}
  {\color{accent}\rule{2.2em}{3pt}\\[6pt]\normalsize\MakeUppercase{\chaptertitlename}~\thechapter}{4pt}{\huge\bfseries}
\titleformat{name=\chapter,numberless}[display]{\headfont\color{ink}}{\color{accent}\rule{2.2em}{3pt}}{4pt}{\huge\bfseries}
\titlespacing*{\chapter}{0pt}{-10pt}{22pt}
\titleformat{\section}{\headfont\Large\bfseries\color{ink}}{\thesection}{0.8em}{}
\titleformat{\subsection}{\headfont\large\bfseries\color{ink}}{\thesubsection}{0.7em}{}
\titleformat{\subsubsection}{\headfont\normalsize\bfseries\color{ink2}}{\thesubsubsection}{0.6em}{}
\titleformat{\paragraph}[runin]{\headfont\normalsize\bfseries\color{ink}}{}{0pt}{}[.]
\titlespacing*{\paragraph}{0pt}{2.2ex plus 0.8ex minus 0.3ex}{0.7em}
\setcounter{secnumdepth}{3}
\setcounter{tocdepth}{2}
% Table of contents: room for three-digit page numbers and for numbers such as 5.10.10
\makeatletter
\renewcommand{\@pnumwidth}{2.1em}
\renewcommand{\@tocrmarg}{3.2em}
\renewcommand*{\l@section}{\@dottedtocline{1}{1.5em}{2.6em}}
\renewcommand*{\l@subsection}{\@dottedtocline{2}{4.1em}{3.6em}}
\makeatother
% No paragraph skip between the entries of the table of contents (the text's \parskip would spread it)
\AtBeginDocument{\addtocontents{toc}{\protect\setlength{\parskip}{0pt}}}

% Running heads
\usepackage{fancyhdr}
\pagestyle{fancy}
\fancyhf{}
\fancyhead[L]{\footnotesize\color{muted}\nouppercase{\leftmark}}
\fancyhead[R]{\footnotesize\color{muted}\iffr{Rapport technique}{Technical report}}
\fancyfoot[L]{\raisebox{-0.25\height}{\includegraphics[height=9pt]{logo-deraison}}\hspace{4pt}{\headfont\bfseries\footnotesize\color{ink}Deraison.ai}}
\fancyfoot[R]{\footnotesize\color{ink2}\thepage}
\renewcommand{\headrulewidth}{0pt}
\fancypagestyle{plain}{\fancyhf{}\fancyfoot[L]{\raisebox{-0.25\height}{\includegraphics[height=9pt]{logo-deraison}}\hspace{4pt}{\headfont\bfseries\footnotesize\color{ink}Deraison.ai}}\fancyfoot[R]{\footnotesize\color{ink2}\thepage}}

\setlength{\parskip}{0.45em}
\setlength{\parindent}{0pt}
\setlength{\emergencystretch}{3em}
\widowpenalty=10000
\clubpenalty=10000

\usepackage[backend=biber,style=authoryear-comp,maxcitenames=2,uniquename=false,uniquelist=false,dashed=false,
  sorting=nyt,giveninits=true]{biblatex}
\addbibresource{refs.bib}
% No small capitals for family names: Roboto Serif has none
\renewcommand*{\mkbibnamefamily}[1]{#1}
\DefineBibliographyExtras{french}{\renewcommand*{\mkbibnamefamily}[1]{#1}}

\usepackage{xurl}
\usepackage[hidelinks]{hyperref}
\hypersetup{colorlinks=true,linkcolor=ink,citecolor=link,urlcolor=link,
  pdfauthor={Warith Harchaoui},pdflang={\reportlang}}
\ifx\reportlang\langfr \usepackage[french,nameinlink,noabbrev]{cleveref}\else \usepackage[nameinlink,noabbrev]{cleveref}\fi

% House helpers
\newcommand{\repo}{\href{https://github.com/warith-harchaoui/pineau-rl}{\texttt{github.com/warith-harchaoui/pineau-rl}}}
% A file or code path: may break after a slash or an underscore (never with a hyphen), so that long paths do not
% stick out into the margin
\ExplSyntaxOn
\tl_new:N \l__report_file_tl
\NewDocumentCommand{\file}{m}{
  \tl_set:Nn \l__report_file_tl {#1}
  \tl_replace_all:Nnn \l__report_file_tl {/} {/\allowbreak}
  \tl_replace_all:Nnn \l__report_file_tl {\_} {\_\allowbreak}
  \texttt{\l__report_file_tl}
}
\ExplSyntaxOff
\newcommand{\algo}[1]{\textsf{#1}}
\newcommand{\task}[1]{\textsf{#1}}
% A figure from figures/<name>-<lang>.pdf, full text width, at most \figmaxheight tall (tall diagrams shrink to keep
% room for their caption on the page instead of running into the footer)
\newlength{\figmaxheight}\AtBeginDocument{\setlength{\figmaxheight}{0.86\textheight}}
\newcommand{\fig}[3][\linewidth]{\begin{figure}[tbp]\centering\includegraphics[width=#1,height=\figmaxheight,keepaspectratio]{#2-\reportlang}\caption{#3}\label{fig:#2}\end{figure}}
% A generated table from gen/<lang>/<name>.tex (the file holds its own caption and label tab:<name>)
\newcommand{\gentable}[1]{\input{gen/\reportlang/#1}}
.
% LuaLaTeX + biber. Fonts: Latin Modern for text, headings and code, Latin Modern Math for formulas;
% the figures are drawn in the same Latin Modern.
\documentclass[11pt,a4paper,oneside]{report}
\usepackage[a4paper,top=2.6cm,bottom=2.8cm,inner=2.7cm,outer=2.7cm,headheight=14pt]{geometry}

% The LaTeX fonts throughout: Latin Modern for text, headings and code, Latin Modern Math for formulas
\usepackage{fontspec}
% The ffi and ffl ligatures are set as ff + i and ff + l: Ghostscript (build.sh) loses their ToUnicode entry,
% and pdftotext, search and copy-paste then read « ofÏciel ». The ff, fi and fl ligatures survive it.
\directlua{fonts.handlers.otf.addfeature{name="noffi",type="multiple",data={["f_f_i"]={"f_f","i"},["f_f_l"]={"f_f","l"}}}}
\defaultfontfeatures{RawFeature=+noffi}
\setmainfont{Latin Modern Roman}
\setsansfont{Latin Modern Sans}
\setmonofont{Latin Modern Mono}
\newcommand{\headfont}{\rmfamily}

\def\langfr{fr}
\newcommand{\iffr}[2]{\ifx\reportlang\langfr #1\else #2\fi}
\ifx\reportlang\langfr \usepackage[french]{babel}\frenchsetup{StandardItemLabels=true,ItemLabels=\textbullet}\else \usepackage[british]{babel}\fi
\usepackage{csquotes}
\usepackage{microtype}
\usepackage{amsmath,amssymb}
\usepackage{unicode-math}
\setmathfont{Latin Modern Math}
\usepackage{siunitx}
\ifx\reportlang\langfr \sisetup{locale=FR,group-minimum-digits=4,mode=text}\else \sisetup{locale=UK,group-separator={,},group-minimum-digits=4,mode=text}\fi

% Colours of the figures (article/figures/draw.py)
\usepackage[table]{xcolor}
\definecolor{ink}{HTML}{15171C}
\definecolor{ink2}{HTML}{4A4F58}
\definecolor{muted}{HTML}{868C96}
\definecolor{rulec}{HTML}{C9CDD3}
\definecolor{panel}{HTML}{F5F6F8}
\definecolor{accent}{HTML}{D6453D}
\definecolor{okgreen}{HTML}{2F9E5B}
\definecolor{sacblue}{HTML}{2F6FD6}
\definecolor{crossqorange}{HTML}{EE8320}
\definecolor{ppotorch}{HTML}{C23B75}
\definecolor{ppojax}{HTML}{109C86}
\definecolor{link}{HTML}{2F6FD6}

\usepackage{graphicx}
\graphicspath{{figures/}}
\usepackage{booktabs,longtable,array,tabularx,multirow,threeparttable}
\usepackage{makecell}
\renewcommand{\arraystretch}{1.15}
\usepackage[font={small},labelfont={bf},labelsep=period,justification=raggedright,singlelinecheck=false]{caption}
% Captions above tables (the generated tables put them first), below figures
\captionsetup[table]{position=top,skip=6pt}
\captionsetup[figure]{position=bottom,skip=8pt}
% Floats: fill text pages before making a float page, never a float page with one small figure and a blank half
\usepackage{placeins}
\usepackage{flafter}% a float never comes before the text that calls it
% A long table never starts with its caption and header alone at the foot of a page
\usepackage{needspace}
\BeforeBeginEnvironment{longtable}{\par\Needspace*{14\baselineskip}}
\renewcommand{\topfraction}{0.88}
\renewcommand{\bottomfraction}{0.6}
\renewcommand{\textfraction}{0.1}
\renewcommand{\floatpagefraction}{0.72}
\setcounter{topnumber}{3}
\setcounter{bottomnumber}{2}
\setcounter{totalnumber}{4}
\setlength{\floatsep}{14pt plus 4pt minus 2pt}
\setlength{\textfloatsep}{18pt plus 4pt minus 4pt}
\setlength{\intextsep}{14pt plus 4pt minus 2pt}
\makeatletter
\setlength{\@fptop}{0pt}% float pages fill from the top
\setlength{\@fpsep}{14pt plus 1fil}
\setlength{\@fpbot}{0pt plus 1fil}
\makeatother
\usepackage{subcaption}
\usepackage{enumitem}
\setlist{itemsep=2pt,topsep=4pt}
\usepackage{listings}
\lstset{basicstyle=\ttfamily\small,columns=fullflexible,keepspaces=true,breaklines=true,
  frame=l,framerule=1.2pt,rulecolor=\color{rulec},xleftmargin=1.2em,backgroundcolor=\color{panel},
  showstringspaces=false,commentstyle=\color{muted},keywordstyle=\color{ink}\bfseries}
\usepackage[most]{tcolorbox}
\newtcolorbox{keybox}[1][]{enhanced,colback=panel,colframe=accent,boxrule=0pt,leftrule=3pt,arc=0pt,
  left=10pt,right=10pt,top=6pt,bottom=6pt,fontupper=\small,#1}
\newtcolorbox{notebox}[1][]{enhanced,colback=white,colframe=rulec,boxrule=0.6pt,arc=2pt,
  left=8pt,right=8pt,top=5pt,bottom=5pt,fontupper=\small,#1}

% Headings
\usepackage{titlesec}
\titleformat{\chapter}[display]{\headfont\color{ink}}
  {\color{accent}\rule{2.2em}{3pt}\\[6pt]\normalsize\MakeUppercase{\chaptertitlename}~\thechapter}{4pt}{\huge\bfseries}
\titleformat{name=\chapter,numberless}[display]{\headfont\color{ink}}{\color{accent}\rule{2.2em}{3pt}}{4pt}{\huge\bfseries}
\titlespacing*{\chapter}{0pt}{-10pt}{22pt}
\titleformat{\section}{\headfont\Large\bfseries\color{ink}}{\thesection}{0.8em}{}
\titleformat{\subsection}{\headfont\large\bfseries\color{ink}}{\thesubsection}{0.7em}{}
\titleformat{\subsubsection}{\headfont\normalsize\bfseries\color{ink2}}{\thesubsubsection}{0.6em}{}
\titleformat{\paragraph}[runin]{\headfont\normalsize\bfseries\color{ink}}{}{0pt}{}[.]
\titlespacing*{\paragraph}{0pt}{2.2ex plus 0.8ex minus 0.3ex}{0.7em}
\setcounter{secnumdepth}{3}
\setcounter{tocdepth}{2}
% Table of contents: room for three-digit page numbers and for numbers such as 5.10.10
\makeatletter
\renewcommand{\@pnumwidth}{2.1em}
\renewcommand{\@tocrmarg}{3.2em}
\renewcommand*{\l@section}{\@dottedtocline{1}{1.5em}{2.6em}}
\renewcommand*{\l@subsection}{\@dottedtocline{2}{4.1em}{3.6em}}
\makeatother
% No paragraph skip between the entries of the table of contents (the text's \parskip would spread it)
\AtBeginDocument{\addtocontents{toc}{\protect\setlength{\parskip}{0pt}}}

% Running heads
\usepackage{fancyhdr}
\pagestyle{fancy}
\fancyhf{}
\fancyhead[L]{\footnotesize\color{muted}\nouppercase{\leftmark}}
\fancyhead[R]{\footnotesize\color{muted}\iffr{Rapport technique}{Technical report}}
\fancyfoot[L]{\raisebox{-0.25\height}{\includegraphics[height=9pt]{logo-deraison}}\hspace{4pt}{\headfont\bfseries\footnotesize\color{ink}Deraison.ai}}
\fancyfoot[R]{\footnotesize\color{ink2}\thepage}
\renewcommand{\headrulewidth}{0pt}
\fancypagestyle{plain}{\fancyhf{}\fancyfoot[L]{\raisebox{-0.25\height}{\includegraphics[height=9pt]{logo-deraison}}\hspace{4pt}{\headfont\bfseries\footnotesize\color{ink}Deraison.ai}}\fancyfoot[R]{\footnotesize\color{ink2}\thepage}}

\setlength{\parskip}{0.45em}
\setlength{\parindent}{0pt}
\setlength{\emergencystretch}{3em}
\widowpenalty=10000
\clubpenalty=10000

\usepackage[backend=biber,style=authoryear-comp,maxcitenames=2,uniquename=false,uniquelist=false,dashed=false,
  sorting=nyt,giveninits=true]{biblatex}
\addbibresource{refs.bib}
% No small capitals for family names: Roboto Serif has none
\renewcommand*{\mkbibnamefamily}[1]{#1}
\DefineBibliographyExtras{french}{\renewcommand*{\mkbibnamefamily}[1]{#1}}

\usepackage{xurl}
\usepackage[hidelinks]{hyperref}
\hypersetup{colorlinks=true,linkcolor=ink,citecolor=link,urlcolor=link,
  pdfauthor={Warith Harchaoui},pdflang={\reportlang}}
\ifx\reportlang\langfr \usepackage[french,nameinlink,noabbrev]{cleveref}\else \usepackage[nameinlink,noabbrev]{cleveref}\fi

% House helpers
\newcommand{\repo}{\href{https://github.com/warith-harchaoui/pineau-rl}{\texttt{github.com/warith-harchaoui/pineau-rl}}}
% A file or code path: may break after a slash or an underscore (never with a hyphen), so that long paths do not
% stick out into the margin
\ExplSyntaxOn
\tl_new:N \l__report_file_tl
\NewDocumentCommand{\file}{m}{
  \tl_set:Nn \l__report_file_tl {#1}
  \tl_replace_all:Nnn \l__report_file_tl {/} {/\allowbreak}
  \tl_replace_all:Nnn \l__report_file_tl {\_} {\_\allowbreak}
  \texttt{\l__report_file_tl}
}
\ExplSyntaxOff
\newcommand{\algo}[1]{\textsf{#1}}
\newcommand{\task}[1]{\textsf{#1}}
% A figure from figures/<name>-<lang>.pdf, full text width, at most \figmaxheight tall (tall diagrams shrink to keep
% room for their caption on the page instead of running into the footer)
\newlength{\figmaxheight}\AtBeginDocument{\setlength{\figmaxheight}{0.86\textheight}}
\newcommand{\fig}[3][\linewidth]{\begin{figure}[tbp]\centering\includegraphics[width=#1,height=\figmaxheight,keepaspectratio]{#2-\reportlang}\caption{#3}\label{fig:#2}\end{figure}}
% A generated table from gen/<lang>/<name>.tex (the file holds its own caption and label tab:<name>)
\newcommand{\gentable}[1]{\input{gen/\reportlang/#1}}
.
% LuaLaTeX + biber. Fonts: Latin Modern for text, headings and code, Latin Modern Math for formulas;
% the figures are drawn in the same Latin Modern.
\documentclass[11pt,a4paper,oneside]{report}
\usepackage[a4paper,top=2.6cm,bottom=2.8cm,inner=2.7cm,outer=2.7cm,headheight=14pt]{geometry}

% The LaTeX fonts throughout: Latin Modern for text, headings and code, Latin Modern Math for formulas
\usepackage{fontspec}
% The ffi and ffl ligatures are set as ff + i and ff + l: Ghostscript (build.sh) loses their ToUnicode entry,
% and pdftotext, search and copy-paste then read « ofÏciel ». The ff, fi and fl ligatures survive it.
\directlua{fonts.handlers.otf.addfeature{name="noffi",type="multiple",data={["f_f_i"]={"f_f","i"},["f_f_l"]={"f_f","l"}}}}
\defaultfontfeatures{RawFeature=+noffi}
\setmainfont{Latin Modern Roman}
\setsansfont{Latin Modern Sans}
\setmonofont{Latin Modern Mono}
\newcommand{\headfont}{\rmfamily}

\def\langfr{fr}
\newcommand{\iffr}[2]{\ifx\reportlang\langfr #1\else #2\fi}
\ifx\reportlang\langfr \usepackage[french]{babel}\frenchsetup{StandardItemLabels=true,ItemLabels=\textbullet}\else \usepackage[british]{babel}\fi
\usepackage{csquotes}
\usepackage{microtype}
\usepackage{amsmath,amssymb}
\usepackage{unicode-math}
\setmathfont{Latin Modern Math}
\usepackage{siunitx}
\ifx\reportlang\langfr \sisetup{locale=FR,group-minimum-digits=4,mode=text}\else \sisetup{locale=UK,group-separator={,},group-minimum-digits=4,mode=text}\fi

% Colours of the figures (article/figures/draw.py)
\usepackage[table]{xcolor}
\definecolor{ink}{HTML}{15171C}
\definecolor{ink2}{HTML}{4A4F58}
\definecolor{muted}{HTML}{868C96}
\definecolor{rulec}{HTML}{C9CDD3}
\definecolor{panel}{HTML}{F5F6F8}
\definecolor{accent}{HTML}{D6453D}
\definecolor{okgreen}{HTML}{2F9E5B}
\definecolor{sacblue}{HTML}{2F6FD6}
\definecolor{crossqorange}{HTML}{EE8320}
\definecolor{ppotorch}{HTML}{C23B75}
\definecolor{ppojax}{HTML}{109C86}
\definecolor{link}{HTML}{2F6FD6}

\usepackage{graphicx}
\graphicspath{{figures/}}
\usepackage{booktabs,longtable,array,tabularx,multirow,threeparttable}
\usepackage{makecell}
\renewcommand{\arraystretch}{1.15}
\usepackage[font={small},labelfont={bf},labelsep=period,justification=raggedright,singlelinecheck=false]{caption}
% Captions above tables (the generated tables put them first), below figures
\captionsetup[table]{position=top,skip=6pt}
\captionsetup[figure]{position=bottom,skip=8pt}
% Floats: fill text pages before making a float page, never a float page with one small figure and a blank half
\usepackage{placeins}
\usepackage{flafter}% a float never comes before the text that calls it
% A long table never starts with its caption and header alone at the foot of a page
\usepackage{needspace}
\BeforeBeginEnvironment{longtable}{\par\Needspace*{14\baselineskip}}
\renewcommand{\topfraction}{0.88}
\renewcommand{\bottomfraction}{0.6}
\renewcommand{\textfraction}{0.1}
\renewcommand{\floatpagefraction}{0.72}
\setcounter{topnumber}{3}
\setcounter{bottomnumber}{2}
\setcounter{totalnumber}{4}
\setlength{\floatsep}{14pt plus 4pt minus 2pt}
\setlength{\textfloatsep}{18pt plus 4pt minus 4pt}
\setlength{\intextsep}{14pt plus 4pt minus 2pt}
\makeatletter
\setlength{\@fptop}{0pt}% float pages fill from the top
\setlength{\@fpsep}{14pt plus 1fil}
\setlength{\@fpbot}{0pt plus 1fil}
\makeatother
\usepackage{subcaption}
\usepackage{enumitem}
\setlist{itemsep=2pt,topsep=4pt}
\usepackage{listings}
\lstset{basicstyle=\ttfamily\small,columns=fullflexible,keepspaces=true,breaklines=true,
  frame=l,framerule=1.2pt,rulecolor=\color{rulec},xleftmargin=1.2em,backgroundcolor=\color{panel},
  showstringspaces=false,commentstyle=\color{muted},keywordstyle=\color{ink}\bfseries}
\usepackage[most]{tcolorbox}
\newtcolorbox{keybox}[1][]{enhanced,colback=panel,colframe=accent,boxrule=0pt,leftrule=3pt,arc=0pt,
  left=10pt,right=10pt,top=6pt,bottom=6pt,fontupper=\small,#1}
\newtcolorbox{notebox}[1][]{enhanced,colback=white,colframe=rulec,boxrule=0.6pt,arc=2pt,
  left=8pt,right=8pt,top=5pt,bottom=5pt,fontupper=\small,#1}

% Headings
\usepackage{titlesec}
\titleformat{\chapter}[display]{\headfont\color{ink}}
  {\color{accent}\rule{2.2em}{3pt}\\[6pt]\normalsize\MakeUppercase{\chaptertitlename}~\thechapter}{4pt}{\huge\bfseries}
\titleformat{name=\chapter,numberless}[display]{\headfont\color{ink}}{\color{accent}\rule{2.2em}{3pt}}{4pt}{\huge\bfseries}
\titlespacing*{\chapter}{0pt}{-10pt}{22pt}
\titleformat{\section}{\headfont\Large\bfseries\color{ink}}{\thesection}{0.8em}{}
\titleformat{\subsection}{\headfont\large\bfseries\color{ink}}{\thesubsection}{0.7em}{}
\titleformat{\subsubsection}{\headfont\normalsize\bfseries\color{ink2}}{\thesubsubsection}{0.6em}{}
\titleformat{\paragraph}[runin]{\headfont\normalsize\bfseries\color{ink}}{}{0pt}{}[.]
\titlespacing*{\paragraph}{0pt}{2.2ex plus 0.8ex minus 0.3ex}{0.7em}
\setcounter{secnumdepth}{3}
\setcounter{tocdepth}{2}
% Table of contents: room for three-digit page numbers and for numbers such as 5.10.10
\makeatletter
\renewcommand{\@pnumwidth}{2.1em}
\renewcommand{\@tocrmarg}{3.2em}
\renewcommand*{\l@section}{\@dottedtocline{1}{1.5em}{2.6em}}
\renewcommand*{\l@subsection}{\@dottedtocline{2}{4.1em}{3.6em}}
\makeatother
% No paragraph skip between the entries of the table of contents (the text's \parskip would spread it)
\AtBeginDocument{\addtocontents{toc}{\protect\setlength{\parskip}{0pt}}}

% Running heads
\usepackage{fancyhdr}
\pagestyle{fancy}
\fancyhf{}
\fancyhead[L]{\footnotesize\color{muted}\nouppercase{\leftmark}}
\fancyhead[R]{\footnotesize\color{muted}\iffr{Rapport technique}{Technical report}}
\fancyfoot[L]{\raisebox{-0.25\height}{\includegraphics[height=9pt]{logo-deraison}}\hspace{4pt}{\headfont\bfseries\footnotesize\color{ink}Deraison.ai}}
\fancyfoot[R]{\footnotesize\color{ink2}\thepage}
\renewcommand{\headrulewidth}{0pt}
\fancypagestyle{plain}{\fancyhf{}\fancyfoot[L]{\raisebox{-0.25\height}{\includegraphics[height=9pt]{logo-deraison}}\hspace{4pt}{\headfont\bfseries\footnotesize\color{ink}Deraison.ai}}\fancyfoot[R]{\footnotesize\color{ink2}\thepage}}

\setlength{\parskip}{0.45em}
\setlength{\parindent}{0pt}
\setlength{\emergencystretch}{3em}
\widowpenalty=10000
\clubpenalty=10000

\usepackage[backend=biber,style=authoryear-comp,maxcitenames=2,uniquename=false,uniquelist=false,dashed=false,
  sorting=nyt,giveninits=true]{biblatex}
\addbibresource{refs.bib}
% No small capitals for family names: Roboto Serif has none
\renewcommand*{\mkbibnamefamily}[1]{#1}
\DefineBibliographyExtras{french}{\renewcommand*{\mkbibnamefamily}[1]{#1}}

\usepackage{xurl}
\usepackage[hidelinks]{hyperref}
\hypersetup{colorlinks=true,linkcolor=ink,citecolor=link,urlcolor=link,
  pdfauthor={Warith Harchaoui},pdflang={\reportlang}}
\ifx\reportlang\langfr \usepackage[french,nameinlink,noabbrev]{cleveref}\else \usepackage[nameinlink,noabbrev]{cleveref}\fi

% House helpers
\newcommand{\repo}{\href{https://github.com/warith-harchaoui/pineau-rl}{\texttt{github.com/warith-harchaoui/pineau-rl}}}
% A file or code path: may break after a slash or an underscore (never with a hyphen), so that long paths do not
% stick out into the margin
\ExplSyntaxOn
\tl_new:N \l__report_file_tl
\NewDocumentCommand{\file}{m}{
  \tl_set:Nn \l__report_file_tl {#1}
  \tl_replace_all:Nnn \l__report_file_tl {/} {/\allowbreak}
  \tl_replace_all:Nnn \l__report_file_tl {\_} {\_\allowbreak}
  \texttt{\l__report_file_tl}
}
\ExplSyntaxOff
\newcommand{\algo}[1]{\textsf{#1}}
\newcommand{\task}[1]{\textsf{#1}}
% A figure from figures/<name>-<lang>.pdf, full text width, at most \figmaxheight tall (tall diagrams shrink to keep
% room for their caption on the page instead of running into the footer)
\newlength{\figmaxheight}\AtBeginDocument{\setlength{\figmaxheight}{0.86\textheight}}
\newcommand{\fig}[3][\linewidth]{\begin{figure}[tbp]\centering\includegraphics[width=#1,height=\figmaxheight,keepaspectratio]{#2-\reportlang}\caption{#3}\label{fig:#2}\end{figure}}
% A generated table from gen/<lang>/<name>.tex (the file holds its own caption and label tab:<name>)
\newcommand{\gentable}[1]{\input{gen/\reportlang/#1}}
.
% LuaLaTeX + biber. Fonts: Latin Modern for text, headings and code, Latin Modern Math for formulas;
% the figures are drawn in the same Latin Modern.
\documentclass[11pt,a4paper,oneside]{report}
\usepackage[a4paper,top=2.6cm,bottom=2.8cm,inner=2.7cm,outer=2.7cm,headheight=14pt]{geometry}

% The LaTeX fonts throughout: Latin Modern for text, headings and code, Latin Modern Math for formulas
\usepackage{fontspec}
% The ffi and ffl ligatures are set as ff + i and ff + l: Ghostscript (build.sh) loses their ToUnicode entry,
% and pdftotext, search and copy-paste then read « ofÏciel ». The ff, fi and fl ligatures survive it.
\directlua{fonts.handlers.otf.addfeature{name="noffi",type="multiple",data={["f_f_i"]={"f_f","i"},["f_f_l"]={"f_f","l"}}}}
\defaultfontfeatures{RawFeature=+noffi}
\setmainfont{Latin Modern Roman}
\setsansfont{Latin Modern Sans}
\setmonofont{Latin Modern Mono}
\newcommand{\headfont}{\rmfamily}

\def\langfr{fr}
\newcommand{\iffr}[2]{\ifx\reportlang\langfr #1\else #2\fi}
\ifx\reportlang\langfr \usepackage[french]{babel}\frenchsetup{StandardItemLabels=true,ItemLabels=\textbullet}\else \usepackage[british]{babel}\fi
\usepackage{csquotes}
\usepackage{microtype}
\usepackage{amsmath,amssymb}
\usepackage{unicode-math}
\setmathfont{Latin Modern Math}
\usepackage{siunitx}
\ifx\reportlang\langfr \sisetup{locale=FR,group-minimum-digits=4,mode=text}\else \sisetup{locale=UK,group-separator={,},group-minimum-digits=4,mode=text}\fi

% Colours of the figures (article/figures/draw.py)
\usepackage[table]{xcolor}
\definecolor{ink}{HTML}{15171C}
\definecolor{ink2}{HTML}{4A4F58}
\definecolor{muted}{HTML}{868C96}
\definecolor{rulec}{HTML}{C9CDD3}
\definecolor{panel}{HTML}{F5F6F8}
\definecolor{accent}{HTML}{D6453D}
\definecolor{okgreen}{HTML}{2F9E5B}
\definecolor{sacblue}{HTML}{2F6FD6}
\definecolor{crossqorange}{HTML}{EE8320}
\definecolor{ppotorch}{HTML}{C23B75}
\definecolor{ppojax}{HTML}{109C86}
\definecolor{link}{HTML}{2F6FD6}

\usepackage{graphicx}
\graphicspath{{figures/}}
\usepackage{booktabs,longtable,array,tabularx,multirow,threeparttable}
\usepackage{makecell}
\renewcommand{\arraystretch}{1.15}
\usepackage[font={small},labelfont={bf},labelsep=period,justification=raggedright,singlelinecheck=false]{caption}
% Captions above tables (the generated tables put them first), below figures
\captionsetup[table]{position=top,skip=6pt}
\captionsetup[figure]{position=bottom,skip=8pt}
% Floats: fill text pages before making a float page, never a float page with one small figure and a blank half
\usepackage{placeins}
\usepackage{flafter}% a float never comes before the text that calls it
% A long table never starts with its caption and header alone at the foot of a page
\usepackage{needspace}
\BeforeBeginEnvironment{longtable}{\par\Needspace*{14\baselineskip}}
\renewcommand{\topfraction}{0.88}
\renewcommand{\bottomfraction}{0.6}
\renewcommand{\textfraction}{0.1}
\renewcommand{\floatpagefraction}{0.72}
\setcounter{topnumber}{3}
\setcounter{bottomnumber}{2}
\setcounter{totalnumber}{4}
\setlength{\floatsep}{14pt plus 4pt minus 2pt}
\setlength{\textfloatsep}{18pt plus 4pt minus 4pt}
\setlength{\intextsep}{14pt plus 4pt minus 2pt}
\makeatletter
\setlength{\@fptop}{0pt}% float pages fill from the top
\setlength{\@fpsep}{14pt plus 1fil}
\setlength{\@fpbot}{0pt plus 1fil}
\makeatother
\usepackage{subcaption}
\usepackage{enumitem}
\setlist{itemsep=2pt,topsep=4pt}
\usepackage{listings}
\lstset{basicstyle=\ttfamily\small,columns=fullflexible,keepspaces=true,breaklines=true,
  frame=l,framerule=1.2pt,rulecolor=\color{rulec},xleftmargin=1.2em,backgroundcolor=\color{panel},
  showstringspaces=false,commentstyle=\color{muted},keywordstyle=\color{ink}\bfseries}
\usepackage[most]{tcolorbox}
\newtcolorbox{keybox}[1][]{enhanced,colback=panel,colframe=accent,boxrule=0pt,leftrule=3pt,arc=0pt,
  left=10pt,right=10pt,top=6pt,bottom=6pt,fontupper=\small,#1}
\newtcolorbox{notebox}[1][]{enhanced,colback=white,colframe=rulec,boxrule=0.6pt,arc=2pt,
  left=8pt,right=8pt,top=5pt,bottom=5pt,fontupper=\small,#1}

% Headings
\usepackage{titlesec}
\titleformat{\chapter}[display]{\headfont\color{ink}}
  {\color{accent}\rule{2.2em}{3pt}\\[6pt]\normalsize\MakeUppercase{\chaptertitlename}~\thechapter}{4pt}{\huge\bfseries}
\titleformat{name=\chapter,numberless}[display]{\headfont\color{ink}}{\color{accent}\rule{2.2em}{3pt}}{4pt}{\huge\bfseries}
\titlespacing*{\chapter}{0pt}{-10pt}{22pt}
\titleformat{\section}{\headfont\Large\bfseries\color{ink}}{\thesection}{0.8em}{}
\titleformat{\subsection}{\headfont\large\bfseries\color{ink}}{\thesubsection}{0.7em}{}
\titleformat{\subsubsection}{\headfont\normalsize\bfseries\color{ink2}}{\thesubsubsection}{0.6em}{}
\titleformat{\paragraph}[runin]{\headfont\normalsize\bfseries\color{ink}}{}{0pt}{}[.]
\titlespacing*{\paragraph}{0pt}{2.2ex plus 0.8ex minus 0.3ex}{0.7em}
\setcounter{secnumdepth}{3}
\setcounter{tocdepth}{2}
% Table of contents: room for three-digit page numbers and for numbers such as 5.10.10
\makeatletter
\renewcommand{\@pnumwidth}{2.1em}
\renewcommand{\@tocrmarg}{3.2em}
\renewcommand*{\l@section}{\@dottedtocline{1}{1.5em}{2.6em}}
\renewcommand*{\l@subsection}{\@dottedtocline{2}{4.1em}{3.6em}}
\makeatother
% No paragraph skip between the entries of the table of contents (the text's \parskip would spread it)
\AtBeginDocument{\addtocontents{toc}{\protect\setlength{\parskip}{0pt}}}

% Running heads
\usepackage{fancyhdr}
\pagestyle{fancy}
\fancyhf{}
\fancyhead[L]{\footnotesize\color{muted}\nouppercase{\leftmark}}
\fancyhead[R]{\footnotesize\color{muted}\iffr{Rapport technique}{Technical report}}
\fancyfoot[L]{\raisebox{-0.25\height}{\includegraphics[height=9pt]{logo-deraison}}\hspace{4pt}{\headfont\bfseries\footnotesize\color{ink}Deraison.ai}}
\fancyfoot[R]{\footnotesize\color{ink2}\thepage}
\renewcommand{\headrulewidth}{0pt}
\fancypagestyle{plain}{\fancyhf{}\fancyfoot[L]{\raisebox{-0.25\height}{\includegraphics[height=9pt]{logo-deraison}}\hspace{4pt}{\headfont\bfseries\footnotesize\color{ink}Deraison.ai}}\fancyfoot[R]{\footnotesize\color{ink2}\thepage}}

\setlength{\parskip}{0.45em}
\setlength{\parindent}{0pt}
\setlength{\emergencystretch}{3em}
\widowpenalty=10000
\clubpenalty=10000

\usepackage[backend=biber,style=authoryear-comp,maxcitenames=2,uniquename=false,uniquelist=false,dashed=false,
  sorting=nyt,giveninits=true]{biblatex}
\addbibresource{refs.bib}
% No small capitals for family names: Roboto Serif has none
\renewcommand*{\mkbibnamefamily}[1]{#1}
\DefineBibliographyExtras{french}{\renewcommand*{\mkbibnamefamily}[1]{#1}}

\usepackage{xurl}
\usepackage[hidelinks]{hyperref}
\hypersetup{colorlinks=true,linkcolor=ink,citecolor=link,urlcolor=link,
  pdfauthor={Warith Harchaoui},pdflang={\reportlang}}
\ifx\reportlang\langfr \usepackage[french,nameinlink,noabbrev]{cleveref}\else \usepackage[nameinlink,noabbrev]{cleveref}\fi

% House helpers
\newcommand{\repo}{\href{https://github.com/warith-harchaoui/pineau-rl}{\texttt{github.com/warith-harchaoui/pineau-rl}}}
% A file or code path: may break after a slash or an underscore (never with a hyphen), so that long paths do not
% stick out into the margin
\ExplSyntaxOn
\tl_new:N \l__report_file_tl
\NewDocumentCommand{\file}{m}{
  \tl_set:Nn \l__report_file_tl {#1}
  \tl_replace_all:Nnn \l__report_file_tl {/} {/\allowbreak}
  \tl_replace_all:Nnn \l__report_file_tl {\_} {\_\allowbreak}
  \texttt{\l__report_file_tl}
}
\ExplSyntaxOff
\newcommand{\algo}[1]{\textsf{#1}}
\newcommand{\task}[1]{\textsf{#1}}
% A figure from figures/<name>-<lang>.pdf, full text width, at most \figmaxheight tall (tall diagrams shrink to keep
% room for their caption on the page instead of running into the footer)
\newlength{\figmaxheight}\AtBeginDocument{\setlength{\figmaxheight}{0.86\textheight}}
\newcommand{\fig}[3][\linewidth]{\begin{figure}[tbp]\centering\includegraphics[width=#1,height=\figmaxheight,keepaspectratio]{#2-\reportlang}\caption{#3}\label{fig:#2}\end{figure}}
% A generated table from gen/<lang>/<name>.tex (the file holds its own caption and label tab:<name>)
\newcommand{\gentable}[1]{\input{gen/\reportlang/#1}}
